\subsubsection{Document Classification and Similarity}\label{sec:document-classification-and-similarity}

After seeing how word embeddings can improve \textbf{Information Retrieval (IR)}, another important application is \textbf{Document Classification and Similarity}. The central idea is simple: \hl{if words have semantic embeddings, then entire documents can also be represented using those embeddings.} This allows us to compare documents based on \textbf{their meaning}, rather than only on the exact words they contain.

\highspace
\begin{flushleft}
    \textcolor{Red2}{\faIcon{exclamation-triangle} \textbf{The problem with Bag-of-Words (BoW)}}
\end{flushleft}
Recall the Bag-of-Words representation. Suppose we have three documents:
\begin{itemize}
    \item \textbf{Document $D_{0}$}
    \begin{center}
        ``\emph{The President greets the press in Chicago.}''
    \end{center}
    \item \textbf{Document $D_{1}$}
    \begin{center}
        ``\emph{Obama speaks to the media in Illinois.}''
    \end{center}
    \item \textbf{Document $D_{2}$}
    \begin{center}
        ``\emph{The band gave a concert in Japan.}''
    \end{center}
\end{itemize}
With Bag-of-Words, every word is treated as an independent symbol. Therefore, the vectors of $D_0$ and $D_1$ share very few identical words. Surprisingly, $D_2$ also shares very few identical words. Consequently, Bag-of-Words considers:
\begin{equation*}
    \mathrm{sim}\left(D_0, D_1\right) \approx \mathrm{sim}\left(D_0, D_2\right)
\end{equation*}
Which is clearly \textbf{wrong}. But, \emph{whi is this wrong?} As humans, we immediately understand that ``\emph{Obama}'' is closely related to ``\emph{President}''. Similarly, ``\emph{speaks}'' and ``\emph{greets}'' are semantically related. Likewise, ``\emph{Illinois}'' and ``\emph{Chicago}'' refer to the same geographical area. The meaning is almost identical even though the words are different.

\highspace
\begin{flushleft}
    \textcolor{Green3}{\faIcon{check-circle} \textbf{Word Embeddings solve this}}
\end{flushleft}
Suppose Word2Vec has learned:
\begin{itemize}
    \item \emph{Obama} $\rightarrow$ $\mathbf{w}_{\text{A}}$
    \item \emph{President} $\rightarrow$ $\mathbf{w}_{\text{B}}$
\end{itemize}
Because they often appear in similar contexts, their embeddings satisfy: $A \approx B$. Similarly, we have: \emph{speaks} $\approx$ \emph{greets}, \emph{media} $\approx$ \emph{press}, and \emph{Illinois} $\approx$ \emph{Chicago}. Therefore, we can represent each document as a vector by averaging the embeddings of its words:
\begin{equation}
    \mathbf{d} = \frac{1}{N} \sum_{i=1}^{N} \mathbf{w}_{i}
\end{equation}
Where $N$ is the number of words in the document, and $\mathbf{w}_{i}$ is the embedding of the $i$-th word. The result is \textbf{one vector representing the entire document}. Notice that this is not CBOW. We're simply \textbf{reusing the learned embeddings}.

\highspace
\textcolor{DarkOrange3}{\faIcon{balance-scale} \textbf{Comparing two documents.}} Once every document has a vector representation, we compare them using the cosine similarity:
\begin{equation*}
    \mathrm{sim}\left(D_i, D_j\right) = \frac{\mathbf{d}_i \cdot \mathbf{d}_j}{\left\|\mathbf{d}_i\right\| \cdot \left\|\mathbf{d}_j\right\|} 
\end{equation*}
If the cosine similarity is high, the documents discuss similar topics.

\highspace
\textcolor{Green3}{\faIcon{question-circle} \textbf{Why is this useful?}} This simple representation can already improve many tasks.
\begin{itemize}
    \item[\textcolor{Green3}{\faIcon{check}}] \textcolor{Green3}{\textbf{Document similarity.}} Find documents discussing the same topic.
    \item[\textcolor{Green3}{\faIcon{check}}] \textcolor{Green3}{\textbf{Document classification.}} Instead of classifying sparse Bag-of-Words vectors, a classifier receives dense semantic embeddings. For example, ``\emph{Politics}'', ``\emph{Sports}'', ``\emph{Technology}'', ``\emph{Medicine}'', etc. Documents about politics tend to occupy a similar region of the embedding space, making them easier to classify.
    \item[\textcolor{Green3}{\faIcon{check}}] \textcolor{Green3}{\textbf{Clustering.}}  Documents discussing similar topics naturally group together.
\end{itemize}
The averaging strategy presented here was one of the \textbf{first practical ways} to represent an entire document using Word2Vec. It demonstrated that embeddings learned from a simple language model could be reused as powerful features for downstream NLP tasks such as similarity search and classification. Today, more advanced methods such as \textbf{Doc2Vec}, \textbf{Sentence-BERT}, and modern transformer-based embedding models learn document or sentence embeddings directly, but they all build on the same core idea introduced by Word2Vec: \textbf{represent text in a continuous semantic vector space where meaning corresponds to geometric proximity}.